\chapter{Пространственно-временная кинетическая факторизация}
\label{ch:v8-spacetime-kinetic-factorization-and-gauge-fixing-gate}

\section{Поперечное и продольное подпространства}

Для ненулевого евклидова импульса введём проекторы
\begin{equation}
 P_{\mathrm L}=\frac{pp^T}{p^2},
 \qquad P_{\mathrm T}=I_4-P_{\mathrm L},
 \qquad P_{\mathrm T}P_{\mathrm L}=0.
 \label{eq:v8-kinetic-transverse-longitudinal-projectors}
\end{equation}
В системе координат $p=(1,0,0,0)$ они равны
$P_{\mathrm L}=\operatorname{diag}(1,0,0,0)$ и
$P_{\mathrm T}=\operatorname{diag}(0,1,1,1)$. Все проекторные тождества
проверены точно.

\section{Гессиан до фиксации калибровочной свободы}

Квадратичная часть калибровочного действия факторизуется как
\begin{equation}
 H_0(p)=K_{\text{к}}\otimes p^2P_{\mathrm T}.
 \label{eq:v8-kinetic-ungauged-factorization}
\end{equation}
Поскольку $\rank K_{\text{к}}=12$ и $\rank P_{\mathrm T}=3$, получаем
\begin{equation}
 \rank H_0=36,
 \qquad \dim\ker H_0=12.
 \label{eq:v8-kinetic-ungauged-rank-nullity}
\end{equation}
Ядро является продольной калибровочной свободой, а не отсутствием
внутренней метрики.

\section{Ковариантная фиксация и точный обратный оператор}

После добавления ковариантного члена с параметром $\xi>0$ имеем
\begin{equation}
 H_{\xi}(p)=K_{\text{к}}\otimes
 p^2\left(P_{\mathrm T}+\xi^{-1}P_{\mathrm L}\right),
 \qquad \rank H_{\xi}=48.
 \label{eq:v8-kinetic-gauge-fixed-hessian}
\end{equation}
Его обратный оператор факторизуется точно:
\begin{equation}
 H_{\xi}^{-1}(p)=K_{\text{к}}^{-1}\otimes
 \frac1{p^2}\left(P_{\mathrm T}+\xi P_{\mathrm L}\right),
 \qquad H_{\xi}H_{\xi}^{-1}=I_{48}.
 \label{eq:v8-kinetic-factorized-inverse}
\end{equation}
Таким образом, $K_{\text{к}}^{-1}$ действительно является внутренним
множителем обратной кинетической формы, но не всей физической мобильностью.

Поперечное ограничение
\begin{equation}
 P_{\mathrm T}H_{\xi}^{-1}P_{\mathrm T}
 =K_{\text{к}}^{-1}\otimes\frac{P_{\mathrm T}}{p^2}
 \label{eq:v8-kinetic-transverse-inverse}
\end{equation}
не зависит от $\xi$, тогда как продольное ограничение пропорционально
$\xi$. Точная проверка при $\xi=1$ и $\xi=2$ подтверждает совпадение
поперечных частей и различие полных обратных операторов.

\begin{theorem}[Точная кинетическая факторизация]
Нефиксированный калибровочный гессиан имеет ранг $36$ и двенадцатимерное
продольное ядро. После ковариантной фиксации он имеет ранг $48$, а его
обратный оператор точно факторизуется на внутренний тензор
$K_{\text{к}}^{-1}$ и пространственно-временной поперечно-продольный
оператор. Физическая поперечная часть не зависит от параметра фиксации.
\end{theorem}

\begin{nogocriterion}[Внутренний обратный тензор не задаёт абсолютную скорость]
Факторизация выделяет $K_{\text{к}}^{-1}$, но не выбирает импульсный или
временной масштаб, корреляционную функцию среды и общий множитель
диссипативной скорости. Продольная часть полного обратного оператора также
зависит от выбора фиксации и не является наблюдаемой мобильностью.
\end{nogocriterion}

\begin{protocol}
\begin{enumerate}
 \item размерность кинетического носителя: \textbf{$12\cdot4=48$};
 \item ранг нефиксированного гессиана: \textbf{$36$};
 \item продольная размерность ядра: \textbf{$12$};
 \item ранг после фиксации: \textbf{$48$};
 \item остаток обращения: \textbf{$0$};
 \item поперечная часть зависит от $\xi$: \textbf{нет};
 \item продольная часть зависит от $\xi$: \textbf{да};
 \item внутренний множитель обратной формы: \textbf{$K_{\text{к}}^{-1}$};
 \item абсолютная мобильность: \textbf{не выведена};
 \item следующий гейт: \textbf{происхождение поперечной шумовой мобильности
       из корреляций среды}.
\end{enumerate}
\end{protocol}

Машинный сертификат сохранён в
\path{s2t_v8_spacetime_kinetic_factorization_and_gauge_fixing_gate_results.json}.